% Brandon Wu
% DRAFT — "everything" deck for the async lecture. Deliberately over-complete:
% every beat from async_plan.md is here as a frame, in act order, plus a large
% Backup section. Intended for manual reordering and cutting.
% Companion code: ~/exp/async_sml (all transcripts machine-checked there).

% DOCUMENT CLASS AND PACKAGE USE
\documentclass[aspectratio=169]{beamer}

\newif\ifcolorlambda
\colorlambdafalse

\newif\ifuseaux
\useauxfalse
\useauxtrue

\newcommand{\auxColor}{08b1d4}     % the color of note boxes and stuff
\newcommand{\presentColor}{6D4FE8} % the primary color of the slide borders
\newcommand{\bgColor}{ece4f0}      % the color of the background of the slide
\newcommand{\darkBg}{8b98ad}
\newcommand{\lambdaColor}{\auxColor}

\colorlambdatrue

\usepackage{comment}
\usepackage{soul}
\usepackage{listings}
\usepackage{makecell}

\setbeamertemplate{itemize items}[circle]

\usepackage{lectureSlides}
%%%%%%%%%%%%%%%%%%%%%%%%%%%%%%%%%%%%%%%%%| <----- Don't make the title any longer than this
\title{Asynchronous Programming} % TODO
\subtitle{Waiting without stopping the world} % TODO
\date{16 July 2026} % TODO
\author{Brandon Wu} % TODO

\graphicspath{ {./img/}, {../../../img/} }

\definecolor{boxColor}{HTML}{fce1c0}
\definecolor{highlightColor}{HTML}{faf89d}

\usetikzlibrary {shapes.symbols}
\usetikzlibrary {calc}
\tikzset{
  every path/.style={line width=0.25mm},
  box/.style={
    rectangle, draw=black, inner sep=0.4cm, thick,
    minimum height=1cm,
    fill=boxColor,
    font=\large,
  },
  hlbox/.style={
    rectangle, draw=black, inner sep=0.4cm, thick,
    minimum height=1cm,
    fill=highlightColor,
  }
}

\newcommand{\topthing}[2]{
  \begin{minipage}[t][#1][t]{\textwidth}
    \vspace{\fill}
    #2
    \vspace{\fill}
  \end{minipage}
}

% \ifweb is provided by the mkWeb build; define it (false) for direct pdflatex.
% csname trick: a literal \ifweb inside the branch would confuse TeX's skipper.
\ifdefined\ifweb\else\expandafter\newif\csname ifweb\endcsname\fi

\begin{document}

\ifweb
    \renewcommand{\pause}{}
\fi

\setbeamertemplate{itemize items}[circle]

{
\begin{frame}[plain]
    \colorlambdatrue
    \titlepage
\end{frame}
}

% \menti{TODO}

\begin{frame}[fragile]
  \frametitle{Lesson Plan}

  \tableofcontents
\end{frame}

%%%%%%%%%%%%%%%%%%%%%%%%%%%%%%%%%%%%%%%%%%%%%%%%%%%%%%%%%%%%%%%%%%%%%%%%%%%%%%
\sectionSlide{1}{The Wall}
%%%%%%%%%%%%%%%%%%%%%%%%%%%%%%%%%%%%%%%%%%%%%%%%%%%%%%%%%%%%%%%%%%%%%%%%%%%%%%

\begin{frame}[fragile]
  \frametitle{Programs So Far}

  Generally, in 150, we have been interested in writing functions which
  run to completion and compute some value along the way.

  \pause
  \vspace{\fill}

  Such programs are often mathematical in nature--they only interface
  with the outside world insofar as they are asked to compute some value,
  but the program execution otherwise does not really care about
  anything having to do with the outside world\footnote{Not too dissimilar
  from some computer scientists.}.

  \pause
  \vspace{\fill}

  These are nice, but only form a subset of interesting programs. There
  are many programs which are tailor-made to deal with humans!
\end{frame}

\begin{frame}[fragile]
  \frametitle{Interactive Programs}

  Such programs are often \textit{interactive}--they entail taking in some
  input, computing some results by talking to other programs or servers,
  and require interactivity to be preserved while potentially doing
  long computations in the background.

  \pause
  \vspace{\fill}

  Writing such interactive programs requires the ability to do
  \term{asynchronous programming}---programs which can run multiple
  things over a period of time.

  \pause
  \vspace{\fill}

  \defBox{}{A \term{concurrent} program is one which can run multiple
  different operations in an interleaved fashion over a period of time.}
\end{frame}



% DEMO CUE: cd ~/exp/async_sml && sml pingpong-callbacks.sml
% ping 3 / ping 1 (pongs return out of order!) / prompt stays alive / quit
% (this artifact IS the lecture's REPL, verified; best shown after the fix acts)
\begin{frame}[fragile]
  \frametitle{A Program You Cannot Write (Yet)}


  We want to write a REPL which only has one command, \code{ping}, which
  causes the REPL to return \code{pong}.

  \pause
  \vspace{\fill}

  This command will also take in an optional argument, which is how many
  seconds to wait before returning \code{pong}!

  \begin{codeblock}[language={}]
repl> help
commands: ping [seconds] | help | quit
repl> ping
pong! (waited 0 seconds)
repl> ping 3
pong! (waited 3 seconds)
repl> quit
\end{codeblock}
\end{frame}

\begin{frame}[fragile]
  \frametitle{The Obvious Implementation}

  {\small
  \begin{codeblock}
    datatype command = Help | Ping of int

    fun ping n =
      ( `wait` n;
        print ("pong! (waited " ^ Int.toString n ^ " seconds)")
      )

    fun repl () =
      ( print "repl> ";
        case `parse` (`input_line` ()) of
          Help => (print helpText; repl ())
        | Ping n => (ping n; repl ())
      )
\end{codeblock}
  }

  For highlighted code, let's assume that these functions exist
  for the moment and not worry too much about how they work.
\end{frame}

\begin{frame}[fragile]
  \frametitle{What Just Happened?}

  Makes sense. We notice something interesting, though:

  \pause
  \vspace{\fill}

  \begin{codeblock}
    repl> ping 3
\end{codeblock}

  \pause
  \vspace{\fill}

  The \code{repl>} prompt isn't shown until 3 seconds later,
  when the \code{ping 3} command finishes!

  \pause
  \vspace{\fill}

  This means we can't possibly run \code{help} or even another
  \code{ping} command until the first one finishes. In other words,
  commands in our REPL are \term{blocking}--they prevent us from
  doing anything else until they finish.
\end{frame}



  % \begin{center}
  % \emph{Nonblocking host operations create independently progressing work.
  % Promises represent their eventual results. Functional combinators describe
  % dependencies among those results. Async/await is direct-style notation
  % for this composition.}
  % \end{center}

\begin{frame}[fragile]
  \frametitle{Where the Blocking Lives}

  {\small
  \begin{codeblock}
    datatype command = Help | Ping of int

    fun ping n =
      ( &wait n&;
        print ("pong! (waited " ^ Int.toString n ^ " seconds)")
      )

    fun repl () =
      ( print "repl> ";
        case parse (&input_line ()&) of
          Help => (print helpText; repl ())
        | Ping n => (ping n; repl ())
      )
\end{codeblock}
  }

  In particular, both of the above highlighted expressions are
  blocking---they will stall the entire program until the world answers.
\end{frame}




% \begin{frame}[fragile]
%   \frametitle{Wall \#2: Blocking Compute}

%   \begin{codeblock}
% fun fib n = if n < 2 then n else fib (n-1) + fib (n-2)
%   \end{codeblock}

%   \pause
%   \vspace{\fill}

%   \code{fib 43} takes about 3 seconds. During those 3 seconds, nothing else
%   happens --- not because the world is slow, but because \emph{we} are.

%   \pause
%   \vspace{\fill}

%   Two different walls: waiting on the world (I/O), and being busy ourselves
%   (compute). Keep them separate --- they will get separate solutions.
% \end{frame}

\begin{frame}[fragile]
  \frametitle{The Real Problem Is Logical, Not Mechanical}

  Ultimately, the problem originates from the fact that SML is
  \term{run-to-completion}. It is only allowed to run one thing
  (function) at a time, and it cannot start progressing on another
  until the first one finishes.

  \pause
  \vspace{\fill}

  In other words, our call to \code{ping : int -> unit} is a monolith
  --- it returns a unit after a certain amount of time, and nothing is allowed
  to happen until it finishes.

  \begin{codeblock}
    | Ping n => (ping n; repl ())
\end{codeblock}

  \pause
  \vspace{\fill}

  In this case, our \code{ping} call is the problem. It blocks the entire loop,
  when in reality we want to be able to proceed to the \code{repl ()} call,
  and then in \code{n} seconds somehow come back and do the work
  we need.
\end{frame}

\begin{frame}[fragile]
  \frametitle{The Key Revelation: To-Do Lists}

  Our key revelation is this:

  \pause
  \vspace{10pt}

  \keyBox{}{\textbf{Asynchronous programs are synchronous programs which
  have to-do lists.}}

  \pause
  \vspace{\fill}

  This means we can solve the problem by simply adding a to-do list to
  the program. Our idea will be to make it so that \code{ping} can
  yield \textit{immediately}, but also write down somewhere what to
  do later, and fulfill it at a later time.
\end{frame}

\begin{frame}[fragile]
  \frametitle{A First Attempt: Total Cooperation}

  Let's build the dumbest thing that could possibly work. The rule:
  \term{never commit to waiting for anything}\footnote{This is not necessarily
  dating advice.}. Instead: check for work, nap briefly, check again, forever\footnote{This is what a job is like.}.

  \pause
  \vspace{\fill}

  First, \code{ping} stops performing and starts \textit{scheduling}:

  \begin{codeblock}
    val pending : (Time.time * int) list ref = ref []

    fun ping n =
      pending := (`now` () + n, n) :: !pending
\end{codeblock}

  \pause
  \vspace{\fill}

  \code{ping} is now instant --- it does nothing except write a deadline
  on the to-do list. The actual pong-printing must happen somewhere
  else, later.
\end{frame}

\begin{frame}[fragile]
  \frametitle{A First Attempt: Total Cooperation}

  Somebody has to read the list:

  {\small
  \begin{codeblock}
    fun do_pongs () =
      let
        (* split into the ones past due and not *)
        val (due, rest) =
          List.partition (fn (t, _) => t <= `now` ()) (!pending)
      in
        pending := rest;
        List.app (fn (_, n) =>
          print ("pong! (waited " ^ Int.toString n ^ " seconds)")
        ) due
      end
\end{codeblock}}
  \pause
  \vspace{\fill}

  Nothing clever: split off everything whose deadline has passed, and
  do it. Note nobody ``fires'' a timer --- \term{time never calls us;
  we check the time.}
\end{frame}

\begin{frame}[fragile]
  \frametitle{A First Attempt: Total Cooperation}

  \begin{codeblock}
    fun tick () =
      ( do_pongs ();
        if `input_ready` ()
        then case `parse` (`input_line` ()) of
               Help   => print helpText
             | Ping n => ping n
        else ();
        `wait` 1;
        tick () )
\end{codeblock}

  \pause
  \vspace{\fill}

  One new assumed ability: \code{input_ready} asks whether a line is
  ready \textit{without committing to read it}. Input is never awaited ---
  only checked.
\end{frame}

\begin{frame}[fragile]
  \frametitle{A First Attempt: It Works!}

  \begin{codeblock}[language={}]
repl> ping 3
repl> ping 1
pong! (waited 1 seconds)
pong! (waited 3 seconds)
repl> quit
\end{codeblock}

  \pause
  \vspace{\fill}

  This actually totally \textit{works}. We can overlap \code{ping}s,
  our prompt never dies, and pongs arrive in the proper order. This is
  a concurrent program in plain, sequential SML.

  \pause
  \vspace{\fill}

  However, everything happens on the \textit{second} scale. If we type
  \code{help} at $t = 0.4$, we get answered at $t = 1.0$. Not to mention,
  the program keeps spinning forever.
\end{frame}

\begin{frame}[fragile]
  \frametitle{A First Attempt: The Bill}

  \begin{itemize}
    \item \term{Lag is quantized to the tick.} Everything is answered up
      to a second late. A smaller tick means that the CPU spins, and
      a larger tick makes lag grow.
    \pause
    \item \term{\code{tick} is a god-function.} It must enumerate every
      concern in the program; each new feature means editing it.
    \pause
    \item 86{,}400 wakeups a day\footnote{Not to mention 525,600 * 60 in a year!}, mostly to do nothing.
  \end{itemize}

  \pause
  \vspace{\fill}

  And one deeper limit: our list can only remember \textit{pongs}. We
  have no way of deferring arbitrary work, like printing something
  different or printing two things in sequence.

  \pause
  \vspace{\fill}

  Here's an idea: what if we were able to defer arbitrary work to
  be done later? What does this sound like?
\end{frame}

\begin{frame}[fragile]
  \frametitle{Rediscovering CPS}

  We have rediscovered the need for continuation-passing style. Before,
  we used them only to write code in a different light--now, we use them
  to write code which is fundamentally "pausable" at certain points.

  \pause
  \vspace{\fill}

  Before, we assumed only the existence of \code{wait : int -> unit}. Now,
  we assume the existence of its CPS cousin, \code{waitThen: int -> (unit -> unit) -> unit},
  which accepts a continuation to run sometime after the wait is complete.

  \pause
  \vspace{\fill}

  Importantly, \code{waitThen} \textbf{yields immediately,} but we will require
  some other entity to run the continuation for us, sometime after ---
  an entity much like our \code{tick} loop, just with a smarter nap.
\end{frame}

\begin{frame}[fragile]
  \frametitle{ping, with waitThen}

  {\small
  \begin{codeblock}
    datatype command = Help | Ping of int

    fun ping n =
      `waitThen` n `(fn () =>`
        print ("pong! (waited " ^ Int.toString n ^ " seconds)")
      `)`

    fun repl () =
      ( print "repl> ";
        case parse (input_line ()) of
          Help => (print helpText; repl ())
        | Ping n => (ping n; repl ())
      )
\end{codeblock}
  }
\end{frame}

\begin{frame}[fragile]
  \frametitle{One More Bug}

  We have one more problem, which is that there is still one more
  bug left to squash!

  \pause
  \vspace{\fill}

  \begin{codeblock}[language={}]
    repl> ping 3
    repl>
\end{codeblock}

  \pause
  \vspace{\fill}

  After we type \code{ping 3}, the prompt returns instantly, but we
  still do not see the pong!

  \pause
  \vspace{\fill}

  The problem is that the call to \code{input_line} is still blocking,
  meaning that our continuation that prints the \code{pong} can't go
  while our \code{repl} is waiting. We need a second ability: the
  ability to do some other stuff while waiting for someone to type
  something.
\end{frame}

\begin{frame}[fragile]
  \frametitle{The Same Move for Input}

  We will apply a similar technique--instead of blocking on the result of
  \code{input_line}, we will hand it a function to call once that line
  arrives, and then immediately return.

  \pause
  \vspace{\fill}

  We assume the existence of \code{inputLineThen : (string -> unit) -> unit}, which accepts a continuation
  to run when the line (optionally) is available.
\end{frame}

\begin{frame}[fragile]
  \frametitle{A Subtle Issue}

  {\small
  \begin{codeblock}
    fun repl () =
      ( print "repl> ";
        `inputLineThen (fn line =>`
          case parse `line` of
            Help => print helpText
          | Ping n => ping n
        `);`
        `repl ()`
      )
\end{codeblock}
  }

  \pause
  \vspace{\fill}

  Only, be careful here. There's a subtle issue.
\end{frame}

\begin{frame}[fragile]
  \frametitle{Spot the Difference}

  \begin{minipage}{0.49\textwidth}
  {\small
  \begin{codeblock}
    fun repl () =
      ( print "repl> ";
        inputLineThen (fn line =>
        case parse line of
          Help => print helpText
        | Ping n => ping n
        );
        repl ()
      )
\end{codeblock}
  }
  \end{minipage}
  \begin{minipage}{0.49\textwidth}
  {\small
  \begin{codeblock}
    fun repl () =
      ( print "repl> ";
        inputLineThen (fn line =>
          (case parse line of
            Help => print helpText
          | Ping n => ping n);
          repl ()
        )
      )
\end{codeblock}
  }
  \end{minipage}

  \pause
  \vspace{\fill}

  What is the difference between these two implementations?

  \pause
  \vspace{\fill}

  In a world where we assume \code{inputLineThen} to return immediately,
  the first actually loops forever and spam prints \code{"repl>"} forever.
\end{frame}


\begin{frame}[fragile]
  \frametitle{The Fix}

  The lesson, which we just learned the hard way: \term{after a
  registration, the only ``later'' is inside the callback}. So the
  recursive call must move \textit{inside}:

  \begin{codeblock}
    fun repl () =
      ( print "repl> ";
        inputLineThen (fn line =>
          ( case parse line of
              Help => print helpText
            | Ping n => ping n
          ; `repl ()` ))
      )
\end{codeblock}

  \pause
  \vspace{\fill}

  Each arrival of a line handles one command and \textit{re-arms} the
  next read. The REPL's ``loop'' is no longer a loop we run --- it's a
  chain of registrations, each planting the next.
\end{frame}



% \begin{frame}[fragile]
%   \frametitle{One Future, No Arrivals}

%   A sharper way to say it:

%   \pause
%   \vspace{\fill}

%   A sequential SML program has exactly \term{one future}: the rest of the
%   expression being evaluated. Nothing happens that the program didn't just
%   do itself. Nothing ever \emph{arrives}.

%   \pause
%   \vspace{\fill}

%   What we need is a program with \term{several pending futures} --- and some
%   way for the world to wake the right one.
% \end{frame}

% \begin{frame}[fragile]
%   \frametitle{The Twist}

%   You might think: SML is just an academic language; real languages can do
%   this.

%   \pause
%   \vspace{\fill}

%   No. \term{JavaScript stands behind the same wall.} One thread.
%   Run-to-completion. No way to pause a running function. A pure-JS busy loop
%   freezes every browser tab it runs in.

%   \pause
%   \vspace{\fill}

%   Everything the JavaScript world calls ``async programming'' is fallout
%   from this exact wall. We are about to re-derive all of it. Remember this
%   slide when the TypeScript appears at the end.
% \end{frame}

%%%%%%%%%%%%%%%%%%%%%%%%%%%%%%%%%%%%%%%%%%%%%%%%%%%%%%%%%%%%%%%%%%%%%%%%%%%%%%
\sectionSlide{2}{Callbacks and Promises}
%%%%%%%%%%%%%%%%%%%%%%%%%%%%%%%%%%%%%%%%%%%%%%%%%%%%%%%%%%%%%%%%%%%%%%%%%%%%%%

\begin{frame}[fragile]
  \frametitle{Two Suppositions}

  Everything we've developed so far is based on two suppositions:

  \pause
  \vspace{\fill}

  First, we assume that we have the ability to wait and accept input
  "in the background", in a way that will not block our REPL. This is
  exhibited by the following functions:
  \begin{codeblock}
    val waitThen : int -> (unit -> unit) -> unit
    val inputLineThen : (string -> unit) -> unit
\end{codeblock}

  \pause
  \vspace{\fill}

  We also assume the ability of some master process which decides when
  to run our continuations.

  \pause
  \vspace{\fill}

  The former is the ability to run \term{non-blocking operations}, and
  the latter is known as the \term{event loop}.
\end{frame}

\begin{frame}[fragile]
  \frametitle{The Two Registries}

  \begin{codeblock}
    (* two registries: our timers, and one pending read *)
    val timers : (Time.time * (unit -> unit)) list ref = ref []
    val read_job : (string -> unit) option ref = ref NONE

    fun waitThen time f = timers := (`now` () + time, f) :: !timers
    fun inputLineThen k = read_job := SOME k
\end{codeblock}

  \pause
  \vspace{\fill}

  Our two ``magic'' primitives are one line each: they perform nothing,
  they file paperwork. All the actual running happens in one place:
\end{frame}

\begin{frame}[fragile]
  \frametitle{The Event Loop}

  {\small
  \begin{codeblock}
    fun loop () =
      case (`try_input` (), !read_job, !timers) of
          (SOME line, SOME k, _) =>
            (read_job := NONE; k line; loop ())
        | (_, _, (t, f) :: rest) =>
            if `now` () >= t
            then (timers := rest; f (); loop ())
            else (`nap` (); loop ())
        | (_, NONE, []) => ()
        | _ => (`nap` (); loop ())
\end{codeblock}
  }

  \pause
  \vspace{\fill}

  Crucially, \code{nap} is \textit{not} the tick loop's blind
  one-second snooze: it \term{blocks until input is ready or the
  earliest deadline passes, whichever first}. This is the promised
  smarter nap --- we wake exactly when there is something to do. No
  quantized lag, no wakeups for nothing.
\end{frame}

\begin{frame}[fragile]
  \frametitle{The REPL, Shattered}

  {\small
  \begin{codeblock}
    fun ping n =
      waitThen n (fn () =>
        print ("pong! (waited " ^ Int.toString n ^ " seconds)")
      )

    fun repl_step line =
      ( case parse line of
          Help => (print helpText)
        | Ping n => (ping n);
        ask ()
      )

    and ask () =
      ( print "repl> ";
        inputLineThen (fn line => repl_step line)
      )
\end{codeblock}
  }
\end{frame}

\begin{frame}[fragile]
  \frametitle{One Registry per Operation?}

  This works, we need one piece of mutable state for every kind of operation
  we anticipate needing to fulfill later (a list of timers
  for each \code{waitThen}, and a single read job for every pending
  \code{inputLineThen}).

  \pause
  \vspace{\fill}

  If we wanted to add more (HTTP requests, DB reads, etc) we would need to
  add more state and more kinds of callbacks, and more cases into our
  event loop.

  \pause
  \vspace{\fill}

  Is there a generic way we can handle this?
\end{frame}

\begin{frame}[fragile]
  \frametitle{A New Program: Page Similarity}

  Let's dispense with the REPL for a moment. Suppose we're interested in
  a very simple sequence of operations: we want to read two names, fetch
  their Wikipedia pages, and then compute a similarity score between the two.

  \pause
  \vspace{\fill}

  \begin{codeblock}
    val name1 = read_line ()
    val name2 = read_line ()
    val page1 = fetch_page name1
    val page2 = fetch_page name2
    val score = compute_similarity page1 page2
\end{codeblock}

  \pause
  \vspace{\fill}

  This is a totally blocking program, as it currently stands. Let's
  shatter it into its constituent parts, in the hopes of letting other
  parts run between steps.
\end{frame}

\begin{frame}[fragile]
  \frametitle{Callback Hell}

  The easy way to do this is to just arm each operation with a
  continuation, per usual. We'll talk about what to do with the
  continuation later.

  \pause
  \vspace{\fill}

  \begin{codeblock}
    read_line (fn name1 =>
      read_line (fn name2 =>
        fetch_page name1 (fn page1 =>
          fetch_page name2 (fn page2 =>
            compute_similarity page1 page2 score
          )
        )
      )
    )
\end{codeblock}

  \pause
  \vspace{\fill}

  This is commonly what is known in JavaScript as
  \term{callback hell}.
\end{frame}
\begin{frame}[fragile]
  \frametitle{Two Problems}

  The code is gross. The problem is similar to what we saw in the
  CPS lecture--we induce one level of indentation per continuation.

  \pause
  \vspace{\fill}

  Additionally, we have to prescribe the order in which the continuations
  run. In reality, it's not so important that we literally fetch the
  first name's page before the second name's page--we should actually
  be able to do them at the same time.

  \pause
  \vspace{\fill}

  We want a more general mechanism for this idea of having pending
  work to do with a value that we don't yet have. This is the idea of
  a \term{promise}.
\end{frame}

\begin{frame}[fragile]
  \frametitle{Promises}

  \defBox{}{A \term{promise} is a value which represents the result of
  an asynchronous operation that we do not yet have.}

  \pause
  \vspace{\fill}

  In other words, a \code{t Promise.t} is a promise which will
  \textit{eventually} contain a value of type \code{t}.

  \pause
  \vspace{\fill}

  This is the general mechanism we were looking for. We previously
  changed our functions from blocking ones to ones that returned immediately,
  but accepted continuations telling us what to do once we had the value.
  This is encoded via:

  \begin{codeblock}
    (* a direct function *)
    'a -> 'b

    (* an asynchronous function *)
    'a -> 'b Promise.t
\end{codeblock}
\end{frame}

\begin{frame}[fragile]
  \frametitle{thenDo}

  The point is that once we have a promise, we are allowed to
  \term{register} a continuation on it, which is automatically
  queued to run once the promise is fulfilled.

  \pause
  \vspace{\fill}

  So, instead of the callback hell we had before, we can now define
  an infix operator \code{thenDo}:
  \begin{codeblock}
    thenDo : 'a Promise.t * ('a -> 'b Promise.t) -> 'b Promise.t
\end{codeblock}

  \pause
  \vspace{\fill}

  \begin{codeblock}
    (read_line ()) thenDo (fn name1 =>
    (read_line ()) thenDo (fn name2 =>
    (fetch_page name1) thenDo (fn page1 =>
    (fetch_page name2) thenDo (fn page2 =>
    compute_similarity page1 page2
    ))))
\end{codeblock}

  \pause
  \vspace{\fill}

  This isn't much better, actually. But, we can improve things:
\end{frame}

\begin{frame}[fragile]
  \frametitle{both}

  \begin{codeblock}
    read_line () thenDo (fn name1 =>
    read_line () thenDo (fn name2 =>
      both (fetch_page name1, fetch_page name2)
        thenDo (fn (page1, page2) =>
          compute_similarity page1 page2)))
\end{codeblock}

  \pause
  \vspace{\fill}

  We can use a primitive called
  \code{both : 'a Promise.t * 'b Promise.t -> ('a * 'b) Promise.t} to
  combine two promises into one. Essentially, it becomes one promise which
  is fulfilled only when both of the original promises are fulfilled.

  \pause
  \vspace{\fill}

  Importantly, this does not specify which fetch finishes first! Both
  are \textit{started} together, and we only compute the score once both
  pages have arrived. (The two reads stay sequential on purpose ---
  there's only one keyboard.)
\end{frame}

\begin{frame}[fragile]
  \frametitle{The Signature}

  \begin{codeblock}
signature PROMISE =
sig
  type 'a t
  type 'a promise = 'a t

  val new     : unit -> 'a promise
  val resolve : 'a promise * 'a -> unit
  val upon    : 'a promise * ('a -> unit) -> unit

  val thenDo  : 'a promise * ('a -> 'b promise) -> 'b promise
  val both    : 'a promise * 'b promise -> ('a * 'b) promise
end
\end{codeblock}
\end{frame}

\begin{frame}[fragile]
  \frametitle{The Promise REPL}

  Let's first see what we're building toward. Assume our two primitives
  now return promises instead of taking callbacks:

  \begin{codeblock}
    val sleep     : int -> unit promise
    val read_line : unit -> string promise
\end{codeblock}

  \pause
  \vspace{\fill}

  And one more operation: \code{upon (p, k)} registers a plain
  observer --- run \code{k} on the value once it exists. (Like
  \code{thenDo}, except \code{k} returns nothing, and no new promise
  is created.)

  \pause
  \vspace{\fill}

  The contract underneath: \term{when a promise resolves, everything
  waiting on it drops into the loop's inbox of ready work}.

  \pause
  \vspace{\fill}

  Before we see how to build this, let's see how it would be used.
\end{frame}

\begin{frame}[fragile]
  \frametitle{The REPL, Rewritten}

  {\small
  \begin{codeblock}
    fun ping n =
      upon (sleep n, fn () =>
        print ("pong! (waited " ^ Int.toString n ^ " seconds)"))

    fun repl () =
      ( print "repl> ";
        upon (read_line (), fn line =>
          case parse line of
            Help   => (print helpText; repl ())
          | Ping n => (ping n; repl ()))
      )
\end{codeblock}
  }

  \pause

  This is a similar shape to the first \code{repl} we wrote, but every wait
  is now a registration, the recursion lives \textit{inside} the continuation,
  and nothing blocks.\footnote{From here on, the runtime owns the event
  loop: every program implicitly ends by entering it --- just as in
  JavaScript, where the host runs your script to completion, then loops.}
\end{frame}

% \begin{frame}[fragile]
%   \frametitle{One Last Convenience}

%   Having to say the magic words at the end of every program is annoying:

%   \begin{codeblock}
%     val () = (repl (); loop ())
% \end{codeblock}

%   \pause
%   \vspace{\fill}

%   So we fold the engine into the runtime, and give programs one entry
%   point:

%   \begin{codeblock}
%     fun run main = (main (); loop ())

%     val () = run repl
% \end{codeblock}

%   \pause
%   \vspace{\fill}

%   This is exactly what JavaScript does for you: your script \textit{is}
%   the prelude, and the host quietly appends the loop call --- which is
%   why you've never seen it. We end where JS begins: the engine
%   disappears into the runtime, and a program is once again just code
%   that registers things.
% \end{frame}

\sectionSlide{3}{A Promise Trace}

%%%%%%%%%%%%%%%%%%%%%%%%%%%%%%%%%%%%%%%%%%%%%%%%%%%%%%%%%%%%%%%%%%%%%%%%%%%%%%
% Animated trace of the promise REPL: one frame per step, stamped from
% async-trace-body.tex.  \tracestep{line}{notif}{waiting}{inbox}{terminal}
%%%%%%%%%%%%%%%%%%%%%%%%%%%%%%%%%%%%%%%%%%%%%%%%%%%%%%%%%%%%%%%%%%%%%%%%%%%%%%

% code-line marker: highlights the current line when its number equals
% \TraceLine, using the SAME visual as the backtick highlight in
% codeblocks (toolbox's bt@HL@box): a rounded hlYellow node, strut
% height, base-anchored.  Zero footprint; drawn before the line's text
% so the code prints on top.
\newcommand{\M}[1]{\ifnum#1=\TraceLine\relax
  \rlap{\raisebox{0pt}[0pt][0pt]{\tikz[baseline=0pt]{%
    \node[anchor=base west, fill=hlYellow, outer sep=0pt, inner xsep=1pt,
          inner ysep=0pt, rounded corners=3pt,
          minimum height=\ht\strutbox+1pt,
          minimum width=\dimexpr\linewidth-1.5em\relax]
         {\raisebox{1pt}{\strut}\strut};}}}%
\fi}

% a promise, drawn as a chip: rounded light-blue box in the waiting panel
\newcommand{\PBox}[1]{\tikz[baseline=(p.base)]{%
  \node[fill=hlBlue, rounded corners=2pt, inner xsep=3pt,
        inner ysep=1.5pt] (p) {\ttfamily\scriptsize #1};}}

% a ready task in the inbox: same chip, green (blue = waiting, green = ready)
\newcommand{\TBox}[1]{\tikz[baseline=(p.base)]{%
  \node[fill=hlGreen, rounded corners=2pt, inner xsep=3pt,
        inner ysep=1.5pt] (p) {\ttfamily\scriptsize #1};}}

\newcommand{\tracestep}[5]{%
  \def\TraceLine{#1}%
  \def\TraceNotif{#2}%
  \def\TraceWaiting{#3}%
  \def\TraceInbox{#4}%
  \def\TraceTerm{#5}%
  % One step of the promise-REPL trace. Set these before \input-ing:
%   \TraceLine    line number to mark (0 = none)
%   \TraceNotif   system/world notification text
%   \TraceWaiting contents of the waiting-on-the-world box
%   \TraceInbox   contents of the loop's todo inbox
%   \TraceTerm    terminal contents so far
\begin{frame}[fragile]
  \frametitle{Trace: The Promise REPL}

  \begin{columns}[T]
    \begin{column}{0.55\textwidth}
      \begin{codeblock}[escapechar=@, basicstyle=\footnotesize\ttfamily]
@\M{1}@fun ping n =
@\M{2}@  upon (sleep n, fn () =>
@\M{3}@    print ("pong! (waited...)"))
@\M{4}@
@\M{5}@fun repl () =
@\M{6}@  ( print "repl> ";
@\M{7}@    upon (read_line (), fn line =>
@\M{8}@      case parse line of
@\M{9}@        Help => (print helpText;
@\M{10}@                 repl ())
@\M{11}@      | Ping n => (ping n;
@\M{12}@                   repl ())) )
@\M{13}@
@\M{14}@val () = repl ()
\end{codeblock}
    \end{column}
    \begin{column}{0.41\textwidth}
      {\scriptsize
      \colorbox{highlightColor}{\parbox{0.95\linewidth}{%
        \textbf{system:} \textit{\TraceNotif}}}

      \vspace{4pt}
      \fcolorbox{black}{boxColor}{\parbox{0.95\linewidth}{%
        \textbf{waiting on the world}\\[1pt]
        {\ttfamily \TraceWaiting}}}

      \vspace{4pt}
      \fcolorbox{black}{boxColor}{\parbox{0.95\linewidth}{%
        \textbf{inbox (todo)}\\[1pt]
        {\ttfamily \TraceInbox}}}
      }
    \end{column}
  \end{columns}

  \vspace{2pt}
  {\scriptsize
  \fcolorbox{black}{white}{\parbox{0.94\textwidth}{%
    \textbf{terminal}\\[1pt]
    {\ttfamily \TraceTerm}\vspace{1pt}}}}
  \vspace{-6pt}
\end{frame}
%
}

\tracestep{14}
  {repl () runs to completion; the runtime then (implicitly) enters
   its event loop.}
  {(nothing)}
  {(empty)}
  {~}

\tracestep{6}
  {repl runs, start to finish, uninterrupted.}
  {(nothing)}
  {(empty)}
  {repl> }

\tracestep{7}
  {read\_line parks ``resolve P1'' with the reader; upon parks the
   handler on P1. repl RETURNS.}
  {reader $\to$ \PBox{P1 [line handler]}}
  {(empty)}
  {repl> }

\tracestep{0}
  {loop: inbox empty, no timers due --- sleep until input arrives.
   (Nothing of ours is running.)}
  {reader $\to$ \PBox{P1 [line handler]}}
  {(empty)}
  {repl> }

\tracestep{0}
  {WORLD: user types ``ping 2''. Loop wakes; reader fires resolve P1;
   P1's waiters move to the inbox.}
  {(nothing)}
  {\TBox{handler (``ping 2'')}}
  {repl> ping 2}

\tracestep{8}
  {loop runs the handler: parse the line...}
  {(nothing)}
  {(empty)}
  {repl> ping 2}

\tracestep{11}
  {...it's Ping 2, so call ping 2.}
  {(nothing)}
  {(empty)}
  {repl> ping 2}

\tracestep{2}
  {sleep 2 files a timer that will resolve P2; upon parks the pong
   printer on P2.}
  {timer +2s $\to$ \PBox{P2 [print pong]}}
  {(empty)}
  {repl> ping 2}

\tracestep{12}
  {...then repl () re-arms: new prompt, new read (P3). The handler
   RETURNS; the step is over.}
  {timer +2s $\to$ \PBox{P2 [print pong]}\\
   reader\hspace{2.1em} $\to$ \PBox{P3 [line handler]}}
  {(empty)}
  {repl> ping 2\\ repl> }

\tracestep{0}
  {loop: sleep until input OR the +2s deadline, whichever first.
   (Nothing of ours is running.)}
  {timer +2s $\to$ \PBox{P2 [print pong]}\\
   reader\hspace{2.1em} $\to$ \PBox{P3 [line handler]}}
  {(empty)}
  {repl> ping 2\\ repl> }

\tracestep{0}
  {WORLD: 2 seconds pass. Timer due; resolve P2 fires; P2's waiter
   moves to the inbox.}
  {reader $\to$ \PBox{P3 [line handler]}}
  {\TBox{print pong}}
  {repl> ping 2\\ repl> }

\tracestep{3}
  {loop runs the pong printer. Then: inbox empty, one reader parked
   --- sleep until input. The REPL lives on.}
  {reader $\to$ \PBox{P3 [line handler]}}
  {(empty)}
  {repl> ping 2\\ repl> pong! (waited 2 seconds)}

\sectionSlide{4}{Building Promises}

\begin{frame}[fragile]
  \frametitle{Three Drafts of a Cache, Three Diseases}

  Before we build promises, let's see what can go wrong. Raw callbacks come with
  \term{no guarantees}: nothing checks how a function treats the
  continuation it was handed.

  \pause
  \vspace{\fill}

  Here's an innocent task: fetching a page is slow, so cache it ---
  call \code{k} straight from the cache when we can.

  \pause
  \vspace{\fill}

  We'll write it three times, and each time see that it breaks a
  different unwritten rule.
\end{frame}

\begin{frame}[fragile]
  \frametitle{Draft 1: The Dropped Continuation}

  \begin{codeblock}
    fun fetch_cached name k =
      if `valid` name
      then fetch_page name (fn page =>
             (`insert` (name, page);   (* stash for next time *)
              k page))
      else print "bad name!"
\end{codeblock}

  \pause
  \vspace{\fill}

  On a bad name, \code{k} is never called. The caller doesn't crash ---
  it \term{waits forever, silently}. Everything that was going to
  happen next was inside \code{k}, and nobody ran it.

  \pause
  \vspace{\fill}

  Why it's so nasty: the failure is an \textit{absence}. No error, no
  stack trace, no line to point at --- the program didn't die, it just
  stopped having a future.
\end{frame}

\begin{frame}[fragile]
  \frametitle{Draft 2: The Double Call}

  \begin{codeblock}
    fun fetch_cached name k =
      ( case `lookup` name of
          SOME page => k page    (* meant to stop here! *)
        | NONE => ()
      ; fetch_page name (fn page =>
          (`insert` (name, page); k page)) )
\end{codeblock}

  \pause
  \vspace{\fill}

  On a cache hit, \code{k} runs now --- \textit{and again} when the
  fetch lands. But \code{k} is the rest of the program: two pongs, two
  prompts, two of everything downstream.

  \pause
  \vspace{\fill}

  Nothing errors at the call site; the wreckage appears later,
  elsewhere. (In JavaScript this is the famous ``forgot the
  \code{return} before the callback'' --- arguably the most-committed
  callback bug in history.)
\end{frame}

\begin{frame}[fragile]
  \frametitle{Draft 3: Sometimes Now, Sometimes Later}

  Third try --- and this one \textit{looks} correct:

  \begin{codeblock}
    fun fetch_cached name k =
      case `lookup` name of
        SOME page => k page               (* HIT: k runs NOW   *)
      | NONE => fetch_page name (fn page =>
          (`insert` (name, page); k page)) (* MISS: k runs LATER *)
\end{codeblock}

  \pause

  Now look at it from a \textit{caller's} seat:

  \begin{codeblock}
    ( fetch_cached name (fn page => print page)
    ; print "done!" )
\end{codeblock}

  Which prints first, the page or the done? On a \term{hit},
  \code{print page} runs first, but on a miss, it runs second.
  It's all dependent on the cache's mood.
\end{frame}

\begin{frame}[fragile]
  \frametitle{...and One Disease of Our Own}

  The library isn't the only one who can sin --- callers can too.
  Remember our spam-printing repl? Its bug was the code placed
  \textit{after} the registration --- the \code{repl ()} call ran
  immediately. Here is that same disease, minimized:

  \begin{codeblock}
    ( inputLineThen (fn line => print line)
    ; print "done!" )
\end{codeblock}

  \pause
  \vspace{\fill}

  \code{done!} prints \textit{first}. Registration returns instantly;
  the code after it runs \textit{now}, not after the line arrives.

  \pause
  \vspace{\fill}

  \keyBox{}{After a registration, the only ``later'' is inside the
  callback.}
\end{frame}

\begin{frame}[fragile]
  \frametitle{The Tally}

  \begin{itemize}
    \item \term{Dropped}: some branch never calls \code{k} --- the
      workflow waits forever, silently.
    \item \term{Doubled}: some branch calls \code{k} twice --- the rest
      of the program happens twice.
    \item \term{Sometimes-now}: \code{k} runs before the call returns
      on some inputs --- order at the call site depends on hidden
      state.
    \item \term{Premature after}: code placed after a registration runs
      before the result exists.
  \end{itemize}

  \pause
  \vspace{\fill}

  None of these require carelessness --- every \textit{branch} of every
  callback-taking function must maintain ``\term{exactly once, always
  later}'' by hand, forever, in every library you depend on.

  \pause
  \vspace{\fill}

  Hold that phrase. We're about to make it a \textit{datatype}.
\end{frame}

\begin{frame}[fragile]
  \frametitle{Two Kinds of Queue}

  Our conceptual view of our computational model is that we have
  two kinds of queues:

  \pause
  \vspace{\fill}

  \begin{itemize}
    \item The \term{run queue} (one, owned by the event loop):
      contains jobs, which are ready to run anytime.
    \item The \term{wait queues} (one inside every pending promise):
      continuations linked to each promise, which are ready to receive
      the promise's future value. These cannot run yet, as we do not yet have the value
  \end{itemize}

  \pause
  \vspace{\fill}

  Essentially, a promise contains potential run queue functions,
  which are only promoted once the promise is fulfilled.
\end{frame}

\begin{frame}[fragile]
  \frametitle{Implementing the Run Queue}

  Here is how we implement the run queue:

  \begin{codeblock}
    val todo : (unit -> unit) list ref = ref []

    fun defer f = todo := !todo @ [f]
\end{codeblock}

  \pause
  \vspace{\fill}

  This is the generality we've been chasing since the tick loop,
  finally delivered: \term{deferred computation with no strings
  attached}. A \code{todo} entry is simply code to run, later, no
  qualifications.

  \pause
  \vspace{\fill}

  The loop drains it first, before anything else:

  \begin{codeblock}
    fun loop () =
      case !todo of
          f :: fs => (todo := fs; f (); loop ())
        | [] => (* input, timers, nap: as before *) ...
\end{codeblock}

  This is the first and \term{last} change the loop will ever need.
\end{frame}



\begin{frame}[fragile]
  \frametitle{Implementing Promises}

  Promises sound like magic, but we've essentially already built them.

  \pause
  \vspace{\fill}

  \begin{codeblock}
    val read_job : (string -> unit) option ref = ref NONE
\end{codeblock}

  \pause
  \vspace{\fill}

  A slot where a continuation waits for a string is essentially a promise
  with two limitations: it forgets the value once delivered, and it holds
  only one waiter.

  \pause

  \begin{codeblock}
    datatype 'a state =
        Pending of ('a -> unit) list
      | Resolved of 'a

    type 'a promise = 'a state ref
\end{codeblock}
\end{frame}

\begin{frame}[fragile]
  \frametitle{upon and resolve}

  {\small
  \begin{codeblock}
    fun upon (p, k) =
      case !p of
          Resolved v => defer (fn () => k v) (* 1 *)
        | Pending ks => p := Pending (k :: ks)

    fun resolve (p, v) =
      case !p of
          Resolved _ => raise Fail "resolved twice!" (* 2 *)
        | Pending ks =>
            ( p := Resolved v
            ; List.app (fn k => defer (fn () => k v)) ks )
\end{codeblock}
  }

  \pause
  \vspace{\fill}

  This construction carefully avoids the bugs we previously
  outlined. \code{1} ensures we always defer our continuation,
  even if the answer is present, and \code{2} ensures resolving twice
  raises so we do not double-call.
\end{frame}

\begin{frame}[fragile]
  \frametitle{Wrapping the Primitives}

  Where do promises come from? We wrap the primitives we already have:

  \begin{codeblock}
    (* an empty promise is an empty queue of continuations *)
    fun new () = ref (Pending [])

    fun sleep n =
      let val p = new ()
      in waitThen n (fn () => resolve (p, ())); p end

    fun read_line () =
      let val p = new ()
      in inputLineThen (fn line => resolve (p, line)); p end
\end{codeblock}

  \pause
  \vspace{\fill}

  Notice what did \textit{not} change: \code{waitThen},
  \code{inputLineThen}, the timers list, the read job. Promises
  replaced nothing --- they \term{organize} what we had. A promise is a
  parking spot; \code{resolve} is a courier to the same old loop.
\end{frame}

% \begin{frame}[fragile]
%   \frametitle{The Cache, Finally}

%   Remember the cache we couldn't write? The trick: cache the
%   \term{promise}, not the page.

%   \begin{codeblock}
%     val cache : (string * page promise) list ref = ref []

%     fun fetch_cached name =
%       case `lookup` name of
%           SOME p => p          (* a hit -- even mid-flight! *)
%         | NONE =>
%             let val p = fetch_page name
%             in `insert` (name, p); p end
% \end{codeblock}

%   \pause
%   \vspace{\fill}

%   Every disease is now \textit{structurally impossible}: there is no
%   \code{k} to drop or double-call (callers attach their own), and
%   \code{upon} always defers. And something genuinely new: a second
%   request for a page \term{still in flight} just shares the same
%   promise --- one fetch, many askers, late arrivals served from memory.
%   With raw callbacks, that means hand-rolling a waiter list...
%   which is reinventing promises.
% \end{frame}

\begin{frame}[fragile]
  \frametitle{What Changed in the Event Loop?}

  We saw earlier that we changed the inbox to make it more general.

  \pause
  \vspace{\fill}

  The real advantage we have gained here is that promises both
  compose nicely and also prevent us from committing a class of
  bugs (the ones we saw earlier).

  \pause
  \vspace{\fill}

  And nothing about \textit{waiting} changed: the loop still sleeps
  until input or the next deadline --- promises are a library
  \textit{above} the loop, not machinery inside it. Any future kind of
  operation (HTTP requests, database reads, ...) is one
  \code{sleep}-style wrapper away; the loop gains no new cases, ever.
\end{frame}

%%%%%%%%%%%%%%%%%%%%%%%%%%%%%%%%%%%%%%%%%%%%%%%%%%%%%%%%%%%%%%%%%%%%%%%%%%%%%%
% \sectionSlide{5}{Promises}
% %%%%%%%%%%%%%%%%%%%%%%%%%%%%%%%%%%%%%%%%%%%%%%%%%%%%%%%%%%%%%%%%%%%%%%%%%%%%%%

% \begin{frame}[fragile]
%   \frametitle{What a Promise Is (and Is Not)}

%   A promise represents:
%   \begin{itemize}
%     \item \term{one execution}, already started --- not a reusable recipe;
%     \item one eventual result, \term{shared} by all observers;
%     \item a result \term{remembered} after resolution.
%   \end{itemize}

%   \pause
%   \vspace{\fill}

%   Contrast with two things you know:
%   \begin{itemize}
%     \item \code{unit -> 'a} --- a \emph{recipe}: run it again, it computes
%       again;
%     \item \code{'a susp} (lazy lecture) --- \emph{you} pull, when you choose;
%       a promise is \emph{pushed} by the world, and you cannot force it.
%   \end{itemize}

%   \pause
%   \vspace{\fill}

%   % TS shadow, verbal: `const p = fetch(url)` has ALREADY started; awaiting
%   % p twice does not refetch. async is a return-TYPE statement.
% \end{frame}

% % OPTIONAL aside, if a student asks "can't anyone resolve anyone's promise?":
% % real APIs split the read and write ends into two types (same ref, sealed
% % behind a signature) so only the producer may fulfill --- the modules
% % lecture applied to time. jQuery's Deferred had .resolve() on the object
% % and it was abused; standard JS promises moved resolve into a closure
% % (the executor), and ES2024's Promise.withResolvers() returns the pair.

% \begin{frame}[fragile]
%   \frametitle{The Split Is Forced, Not Chosen}

%   Why not one queue?

%   \pause
%   \vspace{\fill}

%   \begin{itemize}
%     \item \term{Types}: \code{defer} runs \code{unit -> unit} thunks. A
%       parked continuation is \code{'a -> unit} --- and only the promise will
%       ever have the \code{'a}. You literally cannot put it on the run queue
%       yet.
%     \pause
%     \item \term{The alternative is busy-polling}: ``resolved yet? re-enqueue
%       me.'' Burns CPU, and the loop never empties --- so you can no longer
%       tell a waiting program from a running one, or detect that the program
%       is finished, or detect deadlock.
%   \end{itemize}

%   \pause
%   \vspace{\fill}

%   Waiting must be a \term{state}, not an activity.
% \end{frame}

% % \begin{frame}[fragile]
% %   \frametitle{You Have Seen This Graph Before}

% %   Scatter the wait queues across all promises and look at what you have:
% %   nodes (promises), edges (parked continuations: ``needs that value'').

% %   \pause
% %   \vspace{\fill}

% %   It is a \term{dependency DAG} --- the work/span graph from the sequences
% %   unit, except now it is a \emph{runtime data structure} you could walk in
% %   the heap.

% %   \pause
% %   \vspace{\fill}

% %   \begin{itemize}
% %     \item The run queue is the DAG's \term{ready frontier}.
% %     \item \code{resolve} advances the frontier.
% %     \item Multi-source: every host op mints a world-fired source. No single
% %       root --- that plurality \emph{is} the concurrency.
% %     \item Not a forest of trees: \term{joins} (\code{both}) give nodes two
% %       parents. Join-free programs are chains --- the case callbacks survive.
% %       Tree-ness fails exactly at the joins, and the joins are why promises
% %       had to exist.
% %   \end{itemize}
% % \end{frame}

% % \begin{frame}[fragile]
% %   \frametitle{Combinators Are Dependency Edges}

% %   \begin{codeblock}
% % infix 1 thenDo thenMap
% % fun p thenDo next = Promise.bind p next
% % fun p thenMap f   = Promise.map f p
% % \end{codeblock}

% %   \pause
% %   \vspace{\fill}

% %   \begin{itemize}
% %     \item \code{p thenDo f} --- the next \emph{async} step needs p's value.
% %     \item \code{p thenMap f} --- only \emph{pure} computation remains after p.
% %     \item \code{both (p, q)} --- joins two \emph{already-started} operations.
% %     \item \code{upon p k} --- observe, without making your workflow depend
% %       on p.
% %   \end{itemize}

% %   \pause
% %   \vspace{\fill}

% %   One warning, now and twice more before the end: \term{combinators create
% %   no concurrency.} The calls that made p and q did. These only describe
% %   dependencies among results.
% % \end{frame}

% \begin{frame}[fragile]
%   \frametitle{The Counter Dance, Retired}

%   \begin{codeblock}
% fun both (p, q) =
%   bind p (fn x => map (fn y => (x, y)) q)
% \end{codeblock}

%   \pause
%   \vspace{\fill}

%   Every hand-rolled join in every 2010 Node codebase --- the two slots, the
%   \code{arrive} check, the off-by-one bugs --- replaced by two lines,
%   written once, in the library.

%   \pause
%   \vspace{\fill}

%   This is the single biggest thing promises bought. Not prettiness:
%   \term{fan-in as a value}.
% \end{frame}

% \begin{frame}[fragile]
%   \frametitle{Honest Accounting}

%   Fixed, permanently:
%   \begin{itemize}
%     \item joins (\code{both}), exactly-once, sharing, memoization;
%     \item the result is a \term{value} again --- \code{main} returns
%       \code{unit Promise.t}; callers can compose it.
%   \end{itemize}

%   \pause
%   \vspace{\fill}

%   Not fixed:
%   \begin{itemize}
%     \item the nesting is still CPS --- ``flat'' promise chains are a
%       formatting idiom, not a structural change;
%     \item \code{NONE} handled per step --- exactly the tax JS's
%       \emph{rejection channel} exists to remove (backup slide).
%   \end{itemize}

%   \pause
%   \vspace{\fill}

%   Fixing the nesting for real is the last act of this lecture.
% \end{frame}

% %%%%%%%%%%%%%%%%%%%%%%%%%%%%%%%%%%%%%%%%%%%%%%%%%%%%%%%%%%%%%%%%%%%%%%%%%%%%%%
% \sectionSlide{6}{The Jobs REPL}
% %%%%%%%%%%%%%%%%%%%%%%%%%%%%%%%%%%%%%%%%%%%%%%%%%%%%%%%%%%%%%%%%%%%%%%%%%%%%%%

% \begin{frame}[fragile]
%   \frametitle{The Public Runtime}

%   Application code sees only this:

%   \begin{codeblock}
% signature ASYNC_RUNTIME =
% sig
%   structure Promise : PROMISE

%   val inputLine : unit -> string option Promise.t
%   val spawn     : (unit -> 'a) -> 'a Promise.t
%   val print     : string -> unit

%   (* drive to quiescence; warn if main never finished *)
%   val run : unit Promise.t -> unit
% end
% \end{codeblock}

%   \pause
%   \vspace{\fill}

%   No threads. No sockets. No scheduler. Two async operations, an output,
%   and a way to start the world turning.
% \end{frame}

% \begin{frame}[fragile]
%   \frametitle{Launching a Job}

%   \begin{codeblock}
% fun launch n =
%   let
%     val id = freshJob ()
%     val () = print ("[" ^ Int.toString id ^ "] started\n")
%     val result = spawn (fn () => compute n)
%   in
%     Promise.upon result (report id)
%   end
% \end{codeblock}

%   \pause
%   \vspace{\fill}

%   Read the dependency structure: \code{spawn} starts the work and returns
%   the handle \emph{immediately}. \code{upon} attaches the report ---
%   \term{without making anything wait for it}. Then \code{launch} returns.
% \end{frame}

% \begin{frame}[fragile]
%   \frametitle{The Input Loop}

%   \begin{codeblock}
% fun loop () =
%   let
%     val () = print "jobs> "
%   in
%     inputLine () thenDo
%       (fn NONE => Promise.pure ()
%         | SOME line => dispatch (parse line))
%   end

% and dispatch (Run n) = (launch n; loop ())
%   | dispatch Help    = (print helpText; loop ())
%   | dispatch Quit    = Promise.pure ()
%   | dispatch Invalid = (print "invalid command\n"; loop ())
% \end{codeblock}

%   \pause

%   One real dependency (\code{thenDo}); one deliberate non-dependency
%   (\code{launch n}).
% \end{frame}

% \begin{frame}[fragile]
%   \frametitle{The Bug That Survives All Abstraction}

%   Now the variant. Same promises, same spawn, same everything --- one line
%   differs:

%   \begin{codeblock}
%   (* responsive *)
% | dispatch (Run n) = (launch n; loop ())

%   (* serial *)
% | dispatchSerial (Run n) =
%     launchAwaited n thenDo (fn () => loopSerial ())
% \end{codeblock}

%   \pause
%   \vspace{\fill}

%   % DEMO CUE: sml live-run-serial.sml -- run 43, then type help, watch it hang
%   The serial version awaits each job before prompting again.
% \end{frame}

% \begin{frame}[fragile]
%   \frametitle{Same Script, Two Fates}

%   Script: \code{run 40} / \code{help} / job 0 completes / \code{quit}

%   \begin{columns}[T]
%   \begin{column}{0.48\textwidth}
%   \begin{codeblock}[language={}]
% jobs> run 40
% [0] started
% jobs> help
% commands: run n | ...
% jobs> [0] finished: 80
% quit
% \end{codeblock}
%   \end{column}
%   \begin{column}{0.48\textwidth}
%   \begin{codeblock}[language={}]
% jobs> run 40
% [0] started
% help
% [0] finished: 80
% jobs> commands: run n | ...
% jobs> quit
% \end{codeblock}
%   \end{column}
%   \end{columns}

%   \pause
%   \vspace{\fill}

%   On the right, \code{help} was typed --- and \emph{ignored} until the job
%   finished. Both transcripts are machine-checked by the same test script.
% \end{frame}

% \begin{frame}[fragile]
%   \frametitle{The Exam Question}

%   \begin{center}
%   \emph{Both versions use promises and independently progressing
%   computation. Why does only one remain responsive?}
%   \end{center}

%   \pause
%   \vspace{\fill}

%   Because the serial version adds an \term{unnecessary dependency edge}:
%   ``read the next command'' now depends on ``current job finishes.''
%   Nothing blocked. The event loop was free the whole time. The
%   \emph{dependency structure} was wrong.

%   \pause
%   \vspace{\fill}

%   \begin{center}
%   \term{await expresses a local data dependency, not global blocking.}
%   \end{center}

%   % TS shadow: this is `for (const u of urls) await fetch(u)` --- the most
%   % common async performance bug in industry. Same structure, same fix.
% \end{frame}

% \begin{frame}[fragile]
%   \frametitle{The Complete State of an Async Program}

%   Freeze the program at any instant. Everything lives in exactly three
%   places:

%   \begin{enumerate}
%     \item the \term{run queue} --- ready thunks (usually empty!);
%     \item the \term{promise wait queues} --- parked continuations, read side;
%     \item the \term{host's registration slots} --- held resolve-callbacks for
%       in-flight operations, write side. Node calls these \emph{handles}.
%   \end{enumerate}

%   \pause
%   \vspace{\fill}

%   Mid-session snapshot of our REPL, two jobs running:

%   \begin{codeblock}[language={}]
% run queue:       []              (idle)
% input promise:   Pending [rest-of-the-REPL]
% job 0 promise:   Pending [report 0]
% job 1 promise:   Pending [report 1]
% host:            reader armed; 2 jobs in flight
% \end{codeblock}

%   ``Doing nothing'' --- yet three futures fully described.
% \end{frame}

% \begin{frame}[fragile]
%   \frametitle{Where Is the REPL Loop Right Now?}

%   Serious question. Point to it. Which data structure holds ``the REPL
%   workflow''?

%   \pause
%   \vspace{\fill}

%   \term{Nowhere. Workflows have no runtime identity.} A workflow exists
%   only as, at most, one parked continuation in some promise's wait queue
%   --- a chain of continuations threaded through the queues over time.

%   \pause
%   \vspace{\fill}

%   Contrast threads: a thread has a stack, an identity, an entry in a
%   scheduler table; you can list threads and kill one.

%   \pause
%   \vspace{\fill}

%   Real-world fallout: Node cannot show you ``all in-flight operations''
%   (only \emph{count} handles); async stack traces are famously bad (the
%   stack dies at every await); ``structured concurrency'' is the movement
%   to give workflows their identity back.
% \end{frame}

% \begin{frame}[fragile]
%   \frametitle{Every Fulfillment Happens On the Loop}

%   Promises never resolve ``in the background.'' Two flavors, both ordinary
%   run-queue steps:

%   \begin{itemize}
%     \item \term{root resolutions} --- the host queues a task whose body
%       \emph{is} the resolve call it has held since \code{fromCallback}. The only
%       doorway for outside values.
%     \item \term{derived resolutions} --- the cascade: \code{map}/\code{bind}
%       adapters resolving downstream promises as their steps run.
%   \end{itemize}

%   \pause
%   \vspace{\fill}

%   Corollary, worth saying aloud: \term{no promise ever changes state
%   underneath a running callback.} The world advances only \emph{between}
%   steps. This is why nothing in today's code needed a lock.
% \end{frame}

% \begin{frame}[fragile]
%   \frametitle{The UI Principle}

%   \begin{codeblock}
% val result = spawn (fn () => compute n)
% val () = Promise.upon result
%            (fn answer => updateUI answer)
% \end{codeblock}

%   \pause
%   \vspace{\fill}

%   Jobs produce \term{values}. Only event-loop callbacks touch the UI /
%   shared state --- and those run one at a time.

%   \pause
%   \vspace{\fill}

%   Invariant you get for free: job \emph{completion order} is
%   nondeterministic, but \emph{UI state transitions} are sequential.
%   Concurrency where it helps, sequentiality where you reason.
% \end{frame}

% \begin{frame}[fragile]
%   \frametitle{An Await Nobody Will Satisfy}

%   What if you await a promise that never resolves? Serial REPL, a job that
%   never completes:

%   \begin{codeblock}[language={}]
% jobs> run 5
% [0] started
% quit
% [runtime] event loop quiescent but main
%           workflow never finished (deadlock?)
% \end{codeblock}

%   \pause
%   \vspace{\fill}

%   The loop goes \term{quiescent}: run queue empty, nothing in flight that
%   could fill it. The parked continuation isn't lost, and it isn't spinning
%   --- it is parked on a condition that can never fire. Because waiting is a
%   \emph{state}, deadlock is \emph{detectable}.
% \end{frame}

%%%%%%%%%%%%%%%%%%%%%%%%%%%%%%%%%%%%%%%%%%%%%%%%%%%%%%%%%%%%%%%%%%%%%%%%%%%%%%
\sectionSlide{5}{Asynchronous Code and TypeScript}
%%%%%%%%%%%%%%%%%%%%%%%%%%%%%%%%%%%%%%%%%%%%%%%%%%%%%%%%%%%%%%%%%%%%%%%%%%%%%%

\begin{frame}[fragile]
  \frametitle{Asynchronous TypeScript}

  Today's lecture involved trying to phrase asynchronous code in
  a run-to-completion single-threaded program, with access to
  higher-order functions.

  \pause
  \vspace{\fill}

  We're lucky that we weren't the first to have to do this--JavaScript
  was one of the first languages to tread this path, long ago!

  \pause
  \vspace{\fill}

  Everything we implemented today has a direct analogue somewhere in
  the history of JavaScript.
\end{frame}


\begin{frame}[fragile]
  \frametitle{async/await, Compiled}


  It turns out that frontend development just ends up being
  functional programming all along.

  \pause
  \vspace{\fill}

  \begin{center}
  \begin{tabular}{lll}
    2009 & callbacks (Node.js) & = CPS \hfill (L11) \\
    2011 & async.js combinators & = HOFs \hfill (L10) \\
    2015 & promises (ES6) & = CPS + single-assignment ref \hfill (L19) \\
    2017 & async/await (ES2017) & = the CPS transform, by compiler \hfill (L11) \\
  \end{tabular}
  \end{center}

  \pause
  \vspace{\fill}

  The very themes you have been learning in this course are the same
  as powered the entire JavaScript and TypeScript ecosystems,
  essentially.
\end{frame}


\begin{frame}[fragile]
  \frametitle{Asynchronous TypeScript}

  TypeScript has a lot of nice abstractions for dealing with
  asynchronous code. You can write:

  \begin{codeblock}
    const result = await fetch("https://example.com");
    const data = await result.json();
    console.log(data);
  \end{codeblock}

  \pause
  \vspace{\fill}

  This code just compiles down to our promise-based equivalent!
  \begin{codeblock}
    fetch "https://example.com" thenDo (fn result =>
      jsonify result thenDo (fn data =>
        print data
      )
    )
  \end{codeblock}
\end{frame}

\begin{frame}[fragile]
  \frametitle{The Reveal}

  All the niceties of asynchronous programming in Typescript, we saw
  today.

  \begin{center}
  \begin{tabular}{ll}
    \textbf{TypeScript} & \textbf{Today} \\
    \texttt{Promise<A>} & \texttt{A Promise.t} \\
    \texttt{p.then(f)} & \texttt{thenMap} or \texttt{thenDo} \\
    \texttt{await p} & \texttt{p thenDo (fn x => ...)} \emph{(compiled!)} \\
    \texttt{Promise.all([p, q])} & \texttt{both (p, q)} \\
    \texttt{async (a: A) => B} & \texttt{A -> B Promise.t} \\
  \end{tabular}
  \end{center}
\end{frame}

% \begin{frame}[fragile]
%   \frametitle{Epilogue: You Were Here First}

%   From Node.js (2009) to async/await (2017), the mainstream spent
%   \term{eight years} climbing out of callback hell, one rung at a time.

%   \pause
%   \vspace{\fill}

%   None of the rungs were new. \term{Promises}: 1988 --- Liskov \&
%   Shrira, for distributed RPC. \term{Direct-style await}: 1994 ---
%   Filinski built it in \textit{SML/NJ}, using first-class
%   continuations. And async/await reached industry through the
%   \term{ML family}: F\# (2007) $\to$ C\# (2012) $\to$ JavaScript,
%   Python, Rust, Swift.

%   \pause
%   \vspace{\fill}

%   Everything you built today was known to this language family before
%   most of you were born. The industry just had to rediscover it ---
%   the hard way.
% \end{frame}

% \begin{frame}[fragile]
%   \frametitle{The Desugaring Rule}

%   \begin{codeblock}[language={}]
% val x = await p            p thenDo (fn x =>
% rest              ==>        rest)
% \end{codeblock}

%   \pause
%   \vspace{\fill}

%   The code after \code{await} \emph{is} the continuation, waiting on the
%   promise.

%   \pause
%   \vspace{\fill}

%   \begin{center}
%   You already knew async/await. It's \code{bind}, wearing syntax.
%   \end{center}
% \end{frame}

% \begin{frame}[fragile]
%   \frametitle{The REPL, in TypeScript}

%   % language=java is close enough for draft highlighting of TS
%   \begin{codeblock}[language=java]
% async function repl(): Promise<void> {
%   while (true) {
%     const line = await inputLine();
%     if (line === null || line === "quit") return;

%     if (isRunCommand(line)) {
%       const job = spawn(() => compute(line));
%       job.then(result => updateUI(result));
%       // deliberately NOT awaited
%     }
%   }
% }
% \end{codeblock}

%   \pause

%   Both moves are here: the genuine dependency (\code{await}), and the
%   deliberate non-dependency (\code{job}). Same structure, new clothes.
% \end{frame}

% \begin{frame}[fragile]
%   \frametitle{But How Does await WORK? Ask the Compiler}

%   JS cannot pause a function. So \code{tsc}, targeting old JavaScript,
%   \emph{rewrites yours}. Real output (typescript@5.5, \code{--target ES5}),
%   outer layer first:

%   \begin{codeblock}[language=java]
% function main() {
%   return __awaiter(this, void 0, void 0, function () {
%     var x, y, _a, a, b;        // locals, hoisted
%     return __generator(this, function (_b) {
%       ...                      // next slide
%     });
%   });
% }
% \end{codeblock}

%   \pause
%   \vspace{\fill}

%   \code{__awaiter} wraps everything in a promise and pumps the machine:
%   run until it yields a promise, \code{.then} the resumption onto it,
%   repeat. Sound familiar? It's \code{fromCallback} plus our driver.
% \end{frame}

% \begin{frame}[fragile]
%   \frametitle{The Machine Itself}

%   \begin{codeblock}[language=java]
% switch (_b.label) {
%   case 0: return [4, inputLine()];
%   case 1:
%     x = _b.sent();
%     return [4, inputLine()];
%   case 2:
%     y = _b.sent();
%     return [4, Promise.all([...])];
%   case 3:
%     _a = _b.sent(); a = _a[0]; b = _a[1];
%     console.log("sum: " + (a + b));
%     return [2];
% }
% \end{codeblock}

%   \code{[4, p]} = ``I yield promise p; resume me at the next label.''
% \end{frame}

% \begin{frame}[fragile]
%   \frametitle{Reading the Machine}

%   \begin{itemize}
%     \item Your body became a \term{switch over numbered states} --- each
%       case is exactly the code between two awaits. The compiler wrote your
%       continuations, then \emph{numbered} them instead of allocating
%       closures.
%     \pause
%     \item \code{await} compiled to \term{return}. Suspension \emph{is}
%       returning --- the only kind of stop JS allows. \code{label} was bumped
%       first, so the machine knows where to resume.
%     \pause
%     \item Locals are \term{hoisted} (\code{var x, y}) so they survive
%       between calls: the stack frame, made manual.
%     \pause
%     \item \code{label} is \term{the function's program counter, evicted to a
%       variable} --- because the function must keep giving the CPU away.
%   \end{itemize}

%   \pause
%   \vspace{\fill}

%   \begin{center}
%   \emph{The TypeScript compiler does to your code what Lecture 11 taught
%   you to do by hand.}
%   \end{center}
% \end{frame}

% \begin{frame}[fragile]
%   \frametitle{The Same Machine, in SML, by Hand (1/2)}

%   \begin{codeblock}
% val x = ref ""  val y = ref ""     (* var x, y  *)
% val label = ref 0                  (* _b.label  *)

% fun step sent =
%   case (!label, sent) of
%       (0, _) =>
%         ( label := 1 ; print "x? "
%         ; upon (map Line (inputLine ())) step )
%     | (1, Line (SOME line)) =>
%         ( x := line ; label := 2 ; print "y? "
%         ; upon (map Line (inputLine ())) step )
%     | ...
% \end{codeblock}

%   \pause

%   Hoisted refs, a label, a case dispatch. Each arm registers
%   \emph{step itself} on the next promise, then falls out of the function.
% \end{frame}

% \begin{frame}[fragile]
%   \frametitle{The Same Machine, in SML, by Hand (2/2)}

%   \begin{codeblock}
%     | (2, Line (SOME line)) =>
%         ( y := line ; label := 3
%         ; upon (map Pair (both ( spawnFib (!x)
%                                , spawnFib (!y) )))
%                step )
%     | (3, Pair (a, b)) =>
%         print ("sum: " ^ Int.toString (a + b) ^ "\n")
%     | _ => raise Fail "impossible state"
% \end{codeblock}

%   \pause

%   This is not an analogy. It is the same machine. (Transcript-identical
%   with all the other stagings --- machine-checked.)
% \end{frame}

% \begin{frame}[fragile]
%   \frametitle{What the Types Expose}

%   One divergence from the JS output, and it's telling:

%   \begin{codeblock}
% datatype sent =            (* _b.sent() : any *)
%     Line of string option
%   | Pair of int * int
% \end{codeblock}

%   \pause
%   \vspace{\fill}

%   JS's \code{sent()} is untyped, so the machine quietly awaits a string
%   here and a pair there. SML makes you declare \term{the union of
%   everything this function will ever await}.

%   \pause
%   \vspace{\fill}

%   That's the classic hidden cost of defunctionalization, surfacing in the
%   types. The compiler knew this union all along --- dynamic typing just let
%   it go unspoken.
% \end{frame}


\begin{frame}[fragile]
  \frametitle{What We Hid (On Purpose)}

  The actual system layer underneath consists of threads, processes,
  sockets, syscalls, none of which we needed to get into to reason
  about the application layer correctly. However, there are a few things
  we glossed over that could improve the design:

  \pause

  \begin{itemize}
    \item \term{Errors}: our promises only succeed. Real ones carry a
      rejection channel (\code{Rejected of exn}), so failures flow down
      chains to one handler --- JS's \code{.catch}.
    \item \term{Type safety}: a proper promise library would separate the
      type of the promise from its resolver, so readers can't resolve.
    \item \term{Compute-heavy jobs}: all our work \textit{waits}. A job
      that \textit{computes} for 3s freezes the loop; fixing that needs
      actual concurrency --- threads, processes, workers.
    \item \term{Cancellation}: nothing can un-register a pending
      workflow. Done right, suspensions become values a runner can
      decline to continue.
    \item \term{Bytes, not lines}: input arrives in chunks, and
      ``readable'' doesn't mean a whole line is there --- real hosts
      keep buffers.
  \end{itemize}
\end{frame}

\begin{frame}[fragile]
  \frametitle{The Whole Lecture, in Five Lines}

  \begin{itemize}
    \item Asynchronous programs are synchronous programs with
      \term{to-do lists}. A main loop, some registries, and
      some well-informed naps.
    \pause
    \item After a registration, the only ``later'' is \term{inside the
      callback} --- code placed after it runs \textit{now}.
    \pause
    \item A \term{promise} is the answer as a first-class value: a list of
      parties waiting, or the answer itself, which is delivered exactly
      once by construction.
    \pause
    \item \code{await} is a \term{compiled promise chain}. It's not
      magic, it's just CPS where you don't have to write the
      continuations.
    \pause
    \item Asynchronous programming boils down to functional programming:
      first-class functions to write the future down, CPS to shatter it,
      and a ref to remember the details.
  \end{itemize}
\end{frame}

% %%%%%%%%%%%%%%%%%%%%%%%%%%%%%%%%%%%%%%%%%%%%%%%%%%%%%%%%%%%%%%%%%%%%%%%%%%%%%%
% \sectionSlide{11}{Backup}
% %%%%%%%%%%%%%%%%%%%%%%%%%%%%%%%%%%%%%%%%%%%%%%%%%%%%%%%%%%%%%%%%%%%%%%%%%%%%%%

% \begin{frame}[fragile]
%   \frametitle{Backup: What If the Computation Raises?}

%   Our promises are success-only. The real model adds a second channel:

%   \begin{codeblock}
% datatype 'a state =
%     Pending of ('a -> unit) list * (exn -> unit) list
%   | Resolved of 'a
%   | Rejected of exn
% \end{codeblock}

%   \pause
%   \vspace{\fill}

%   Rejection is what removes the per-step \code{NONE}/error handling tax you
%   saw in every stage: failures propagate down the chain to one handler ---
%   \code{.catch} --- the way exceptions propagate up a stack. Out of scope
%   today; the design falls out of everything you already know.
% \end{frame}

% % \begin{frame}[fragile]
% %   \frametitle{Backup: Concurrency Is Not Parallelism}

% %   Measured, on the lecture machine: \code{fib 43} alone $\approx$ 2.8s.
% %   Two overlapped \code{fib 43}s: both finish $\approx$ 11s. \term{No
% %   speedup.}

% %   \pause
% %   \vspace{\fill}

% %   Our host interleaves everything on one core. What you saw today is
% %   \term{responsiveness}, never throughput.

% %   \pause
% %   \vspace{\fill}

% %   You already have the parallelism story --- work/span, sequences: the
% %   \emph{deterministic cost} of a known DAG. Today was the concurrency
% %   story: the \emph{nondeterministic unfolding} of one. Keep the axes
% %   separate. (Node is the same: parallel JS does not exist on the main
% %   thread; workers live in the host --- exactly where our parallel host
% %   would go.)
% % \end{frame}

% \begin{frame}[fragile]
%   \frametitle{Backup: The Purity Rule Has Teeth}

%   Because spawned thunks are preemptively interleaved, a thunk that
%   mutates shared refs is \term{observably racy}. Event-loop callbacks keep
%   run-to-completion atomicity; spawned thunks do not.

%   \pause
%   \vspace{\fill}

%   So ``jobs are pure functions that produce values; only loop callbacks
%   touch shared state'' is not a style preference. It is the safety
%   condition under which today's abstraction never leaks.

%   \pause
%   \vspace{\fill}

%   JS enforces the same rule by construction: workers share nothing ---
%   values cross by structured clone. We state it as a contract instead.
% \end{frame}

% \begin{frame}[fragile]
%   \frametitle{Backup: A Genuinely Parallel Host}

%   The host contract never mentions cores --- so a parallel host is a
%   drop-in: fork a process per job, return the result over a pipe. Real
%   multicore speedup; zero changes to any code you saw today.

%   \pause
%   \vspace{\fill}

%   The teachable wrinkle is in the type: a polymorphic \code{'a} cannot
%   cross an address space, so a forking host operation must be monomorphic
%   or take encode/decode functions.

%   \pause
%   \vspace{\fill}

%   Serialization surfacing in the type --- which is exactly what JS's
%   structured-clone requirement for workers is.
% \end{frame}

% \begin{frame}[fragile]
%   \frametitle{Backup: Timers}

%   One host operation we didn't build --- JS's most famous one:

%   \begin{codeblock}
% val startTimer : int * (unit -> unit) -> unit
% val sleep : int -> unit Promise.t
% \end{codeblock}

%   \pause
%   \vspace{\fill}

%   With \code{sleep} and \code{spawn} you can fake \code{fetch}, and then
%   \code{Promise.race}, timeout wrappers, and debounce all port one-to-one
%   from every JS tutorial ever written. (Trivial in the live host; the
%   scripted host would need a logical clock.)
% \end{frame}

% \begin{frame}[fragile]
%   \frametitle{Backup: The Tail-Position Trap}

%   What does this print, and when?

%   \begin{codeblock}
% ( Promise.upon p step
% ; print "done!\n" )
% \end{codeblock}

%   \pause
%   \vspace{\fill}

%   \code{"done!"} prints \emph{immediately} --- while \code{p} is still
%   pending. Registration must be the activation's \term{final act};
%   everything ``after the await'' must live \emph{inside} the handed-off
%   continuation.

%   \pause
%   \vspace{\fill}

%   The compiled form enforces this structurally: \code{return [4, p]}
%   registers and returns in one statement. Hand-written CPS makes it a
%   discipline --- and JS programmers break it weekly with code after
%   \code{.then(...)}.
% \end{frame}

% \begin{frame}[fragile]
%   \frametitle{Backup: A Reified Suspension (the Generator Protocol)}

%   Between promise chains and the state machine sits one more
%   representation --- suspension as an \emph{inspectable value}:

%   \begin{codeblock}
% datatype 'r paused =
%     Val of 'r
%   | Paused of 'r paused Promise.t

% fun drive (Val r)     = deliver r
%   | drive (Paused pp) = Promise.upon pp drive
% \end{codeblock}

%   \pause
%   \vspace{\fill}

%   This is the \term{generator protocol}: \code{Val} = \code{done: true},
%   \code{Paused} = a yielded promise, \code{drive} = the pump. It computes
%   nothing new (the driver's only move is \code{upon} --- it \emph{is}
%   \code{bind}), but a paused workflow is an \emph{inert value}: a driver
%   can log, time-out, interleave, or \emph{decline to continue} ---
%   which is cancellation. And held in your hand, doing nothing, it is the
%   cleanest proof that shattering manufactures no asynchrony.
% \end{frame}

\end{document}
